\documentclass{article}
\usepackage[T1]{fontenc}
\usepackage{lmodern}
\usepackage{graphicx}
\usepackage{amsmath,amssymb}
\usepackage[a4paper,margin=2cm]{geometry}
\usepackage[absolute,overlay]{textpos}
\setlength{\TPHorizModule}{1mm}
\setlength{\TPVertModule}{1mm}
\usepackage[indonesian]{babel}
\usepackage{hyperref}
\setlength{\emergencystretch}{2em}
% unit: Halaman 42

\begin{document}

% === Halaman 42 (llm) ===

\section*{Algoritma Elliptic Curve Diffie-Hellman (ECDH)}

\begin{itemize}
    \item Alice dan Bob ingin berbagi sebuah kunci rahasia $K$ yang sama.
    \begin{itemize}
        \item Alice dan Bob menyepakati kurva eliptik $y^2 = x^3 + ax + b \pmod{p}$ dan titik $B$ pada kurva $B$
        \item Alice dan Bob menghitung kunci publik dan kunci privat masing-masing.
        \begin{itemize}
            \item Alice
            \begin{itemize}
                \item Kunci privat $= a$
                \item Kunci publik $= P_A = a \cdot B$
            \end{itemize}
            \item Bob
            \begin{itemize}
                \item Kunci privat $= b$
                \item Kunci publik $= P_B = b \cdot B$
            \end{itemize}
        \end{itemize}
    \end{itemize}
    \item Alice dan Bob saling mengirim kunci publik masing-masing.
    \item Keduanya melakukan perkalian kunci privatnya dengan kunci publik mitranya untuk mendapatkan kunci rahasia yang mereka bagi
    \begin{itemize}
        \item Alice $\rightarrow K_{AB} = a \cdot P_B = a \cdot (b \cdot B)$
        \item Bob $\rightarrow K_{AB} = b \cdot P_A = b \cdot (a \cdot B)$
        \item Kunci rahasia $= K_{AB} = abB$
    \end{itemize}
\end{itemize}

\vfill
\raggedleft
\small *) Sumber bahan: \textbf{Debdeep Mukhopadhyay}, \textit{Elliptic Curve Cryptography}, \\ Dept of Computer Sc and Engg IIT Madras

\end{document}
